\chapter{From an offer effect to annual contribution}
\label{rewrite:deposit}
\newcommand{\DepositEstimate}{1,183.47}
\newcommand{\DepositLow}{1,128.34}
\newcommand{\DepositHigh}{1,238.60}
\newcommand{\DepositNet}{2.55}
\newcommand{\DepositNetLow}{1.03}
\newcommand{\DepositNetHigh}{4.06}


\section{Fix the intervention before estimating its value}
The preceding accounting calculation defines the question for the first
complete study. Let $Z_i$ denote assignment of the offer to eligible household
$i$, $R_i$ its delivery, $D_i$ acceptance, and $r_i$ the paid annual rate.
These variables describe successive events, not synonyms. Let $B_i(z)$ be
the household's average eligible deposit balance over the following year under
assignment $z$, and let $C_i(z)$ be contribution over that same year. For a
population fixed before assignment, the policy targets are
\[
 \tau_B=\E[B_i(1)-B_i(0)],\qquad
 \tau_C=\E[C_i(1)-C_i(0)].
\]
A statement about the effect among acceptors generally answers neither target.
Acceptance can depend on intentions that also determine balances. Dividing an
assignment effect by an acceptance rate additionally requires an instrumental
variable argument and changes the population to the relevant compliers; the
derivation and restrictions are in Chapter~\ref{r3:instruments}.

Here the constructed design makes a narrower intervention explicit. Every
assigned offer is delivered, and assignment automatically reprices the eligible
balance for one year. Thus $R=Z$ and the price rule is $r=r_0+Z\Delta r$;
acceptance does not gate the payment. A voluntary activation campaign would
need a different expense equation. Assume no interference across households,
no lost outcomes, and unchanged funding margin and direct cost. If the
expected no-offer balance is $B_0$, the incremental expected contribution is
\begin{equation}
 \tau_C=(m-\Delta r)\tau_B-\Delta r B_0-c.
 \label{rewrite:value-map}
\end{equation}
To obtain this expression, subtract $mB_i(0)$ from
$(m-\Delta r)B_i(1)-c$, then take expectations. It charges the rate
increase on the balance that would have remained without the intervention.
The familiar accounting shortcut is justified by this particular price rule.
It would overcharge a campaign that paid only newly acquired funds and would
misdescribe a campaign whose delivery itself failed.

\section{A deliberately transparent assignment design}
The example has two equally weighted customer strata, $X=0,1$. Assignment is
independent of potential balances within each stratum, with probabilities
$e_0=0.5$ and $e_1=0.1$. Both strata have mean no-offer balance
\pounds10,000. The stipulated balance effects are $b_0=0$ and $b_1=2,400$,
with independent normal household noise of standard deviation \pounds400.
These are generated customers. In particular, the independence and normal
law are supplied by construction and have not been inferred from bank data.

Consistency, assignment exchangeability within $X$, and positive assignment
probabilities give
\[
 \tau_B=\sum_{x=0}^1 w_x
   \{\E[B\mid Z=1,X=x]-\E[B\mid Z=0,X=x]\},
 \qquad w_0=w_1=\tfrac12.
\]
The two weights describe the target population, not the observed distribution
of offered customers. The population contrast is therefore \pounds1,200.
A saturated cell-mean estimator uses the same weights. It is a useful first
implementation because every contribution to the answer can be inspected;
with two strata there is no need to learn a complex outcome surface.

A constant residual projection gives a different average when effects vary.
Its weights are proportional to $w_xe_x(1-e_x)$, giving weight
$0.09/(0.25+0.09)$ to the high-benefit stratum and a population coefficient
of about \pounds635.29. That is a valid projection target. Calling it the
equal-stratum population effect would change the financial answer from
positive \pounds3 to about negative \pounds12.53. Chapter~\ref{scisyn:evidence}
retains the separately executed repeated experiments on this target error
and on sparse support. The current single sample demonstrates the full report;
it does not add another calibration experiment.

\section{Carry uncertainty through the expense equation}
Write $n_{xa}$, $\bar B_{xa}$ and $s_{xa}$ for the count, mean and sample
standard deviation of each observed cell. For a total error allowance
$\alpha=0.05$, construct four mean intervals
\[
 I_{xa}=\bar B_{xa}\ \pm\
 t_{n_{xa}-1,1-\alpha/8}\frac{s_{xa}}{\sqrt{n_{xa}}}.
\]
Each interval has conditional noncoverage $\alpha/4$ under its normal
independent sampling law. The union bound implies simultaneous coverage
at least $1-\alpha$, without assuming independence among the four interval
events. Conditioning on the observed assignment counts is legitimate here
because assignment is independent of the outcome errors within strata.
A cell with fewer than two observations returns unsupported rather than a
zero standard error.

If $h_{xa}$ denotes a cell half-width, interval arithmetic yields
\[
 \widehat\tau_B=\tfrac12(\bar B_{01}-\bar B_{00})+
                     \tfrac12(\bar B_{11}-\bar B_{10}),\qquad
 I_B=\widehat\tau_B\ \pm\ \tfrac12\sum_{x,a}h_{xa}.
\]
This construction permits different cell variances. It is conservative and
can become too wide to guide action. The earlier generated study shows why
a narrower pooled interval needs an additional common-variance assumption.
A shared service should expose that tradeoff, not silently choose the
narrower result. The present argument does not justify the Student intervals
under heavy tails, common calendar shocks or dependent households.

With $m=0.03$, $\Delta r=0.0025$, $B_0=10,000$ and $c=5$, equation
\eqref{rewrite:value-map} is increasing in $\tau_B$. Its endpoints
therefore transform directly: $I_C=0.0275 I_B-30$. Uncertainty in these
accounting inputs would need an additional joint set or sensitivity analysis.
Here they are fixed illustrative quantities. The deterministic seed and
8,000 generated observations give the following actual calculation.
\begin{center}
\begin{tabular}{lrr}
\toprule Quantity & Estimate & Interval\\\midrule
Annual average balance effect (\pounds) & \DepositEstimate & [\DepositLow, \DepositHigh]\\
Annual contribution per household (\pounds) & \DepositNet & [\DepositNetLow, \DepositNetHigh]\\
\bottomrule
\end{tabular}
\end{center}
The four cells contain 2,039 unoffered and 1,961 offered households in the
first stratum, and 3,624 unoffered and 376 offered households in the second.
The positive contribution interval supports the action under this declared
model and expense rule. Its margin is small: an additional certain expense
above the lower endpoint removes that interval-based assurance. The result
is unsuitable as a bank forecast; its value is to make the target, support,
uncertainty and accounting inseparable in the completed report.

\section{What the implemented example actually returns}
The local callable estimator takes the three-column generated data, target
weights, four accounting inputs and $\alpha$. It returns cell counts and
uncertainty, the balance target, its interval, contribution and its interval,
and a conditional decision. The same calculation is available at the command
line from the repository root:
\begin{lstlisting}
python3 tools/run_rewrite_examples_20261029.py --report-only
\end{lstlisting}
The callable result and a separate command-line process agree exactly in
the retained execution. The generated data, code, seed and numerical results
are supplied with the study. This is a specific estimator implemented for
the example; it is not evidence that an external causal package was qualified.

The earlier MCP walkthrough establishes graph syntax, conditional
adjustment, information-date checks and refusal of a dependent sampling law.
The public study in Chapter~\ref{gap29:public} additionally sends actual
fitted estimates and covariance summaries through library, CLI and MCP.
It preserves the resulting interval and conditional or contradicted
assignment status. Neither path exposes the deposit raw-data estimator
as an MCP operation.
Chapter~\ref{scisyn:investment} gives the observed requests and the initial
timeout as well as the successful retry. These separate operations do not establish a raw-data-to-financial-report
MCP endpoint.
A future adapter must carry the complete report, including assumptions and
unsupported outcomes, before that claim becomes available.

\section{From the household offer to broader funding evidence}
Two extensions show what this transparent report leaves to be learned.
Chapter~\ref{gap29:public} estimates an assignment contrast from the retained
public experiment, compares adjusted and unadjusted estimators, examines
influence, and subtracts an explicit cost. Its earnings units and population
are different from the deposit account: the exercise tests an inferential
and reporting route, not transport of an earnings effect into bank balances.
Chapter~\ref{gap29:funding} follows institutions rather than households and
asks whether fixed effects can repair a misassigned outcome. They cannot
recover the institution's identity merely by absorbing a persistent level.
The structural funding and calendar arguments explain the additional shock
and dependence assumptions needed for that journey.

For the offer study, the next source of information is assignment support and
correct household aggregation; for the institutional panel it is a coherent
exposure and outcome history. A bank cannot substitute a successful estimate
in one of these populations for the missing design evidence in the other.
The shared report can preserve both questions without pretending they have
the same estimator or sampling law.

\chapter{Servicing decisions change the next decision's information}
\label{rewrite:servicing}

A servicing policy chooses whether to contact a borrower at two dates. Let
$A_0,A_1$ be contact decisions, $L_0$ baseline borrower state and $L_1$ the
state observed between contacts. For this constructed example, $Y$ is
gross contribution in pounds over the two-period horizon, excluding contact
cost. The policy comparison is sustained contact versus no contact, not
the regression coefficient of one contact conditional on every later event.

The exact law developed in Chapter~\ref{r3:longitudinal} is
\begin{align*}
 P(L_0=1)&=0.5,\\
 P(A_0=1\mid L_0)&=0.2+0.6L_0,\\
 P(L_1=1\mid L_0,A_0)&=0.1+0.5L_0+0.2A_0,\\
 P(A_1=1\mid H_1)&=0.15+0.5L_1+0.1A_0,\\
 Y&=20+3A_0+4A_1+10L_1+6L_0.
\end{align*}
Here $H_1=(L_0,A_0,L_1)$ is the history available for the second action.
Independent innovations implement these kernels. The actions have support
for every possible history, and the example has no hidden confounding.
Consequently, replacing both action kernels with a specified policy and
retaining the state-transition law identifies its mean. This is the
g-formula developed from sequential exchangeability in the retained chapter
and its primary treatment. \citep[Chapter 21]{hernan2020}

\section{Average the state produced by each policy}
For a constant policy $a$, the mean intermediate state is
$\E[L_1(a)]=0.1+0.5(0.5)+0.2a=0.35+0.2a$. Hence
\[
 \psi(a)=20+7a+10(0.35+0.2a)+6(0.5)=26.5+9a.
\]
The gross means are \pounds26.50 and \pounds35.50. At a stipulated
\pounds2 per contact, the net contrast is $9-2(2)=\pounds5$.
The reported exact enumeration independently averages the finite state law
and reweights the observed action paths; both recover the same two means.
There is no finite-sample confidence claim in this calculation.

A regression of $Y$ on all four variables correctly recovers action
coefficients 3 and 4. Their sum misses the extra $10(0.2)=2$ contributed
through the change in $L_1$. Removing $L_1$ from the second-stage regression
does not repair the analysis: the historical second action depends on it.
The calculation therefore distinguishes a fitted conditional equation from
the policy mean obtained by averaging that equation over policy-specific
histories. The distinction changes net value from an erroneous \pounds3
to \pounds5 even in this noise-free conditional outcome model.

\section{What changes when the law must be learned}
In an observed study, the kernels and outcome regressions are estimated.
G-computation substitutes fitted transition and outcome laws and averages
under the target policy. Reweighting instead uses ratios of target to
historical action probabilities. Doubly robust constructions combine these
ideas; their longitudinal scores, nuisance restrictions and remainders are
derived in Chapter~\ref{r3:longitudinal}, rather than assumed from their
name. None removes the requirement that the historical records include
the information responsible for action selection.

For sustained contact, a useful diagnostic is the count and weight
distribution of complete histories supporting that policy. Many borrower
rows do not guarantee enough such histories. A policy learned by maximizing
the same estimated value must be evaluated with the additional selection
uncertainty or an honest evaluation sample. An outcome law conditional on
a hidden state introduces another identification problem; a well-fitting
state-space model does not establish that its intervention meaning is correct.

The product implication is to begin with a fixed, short-horizon policy and
observable histories, retaining a direct g-computation and a weighting
comparison. If support is inadequate or an important action determinant is
unrecorded, changing the collection or trial design is more informative than
adding a flexible learner. The positive \pounds5 example explains the
accounting; it does not establish a real contact policy's benefit.

\section{Estimated policies, missing rewards and decision uncertainty}
Chapters~\ref{gap29:servicing} and \ref{gap29:policy-uncertainty} carry the
policy question into fitted data. They use a separate three-decision economy,
with continuous states and rewards already net of the stipulated action cost.
Those reward units are not the pounds in this two-contact example. Whole
customers, rather than isolated rows, supply the sampling units. A separate
training sample chooses the policy; evaluation and resampling hold that
selected policy fixed.

The fitted comparison is informative precisely because it can fail.
Under its maintained state and observation model, improvement coverage is
94\% with an average interval width of 0.711 reward units. Future-dependent
dropout reduces coverage to 6\%; omitting a state predictor reduces it to
61\%. These observed fractions have Monte Carlo uncertainty, reported in
the study. They demonstrate that resampling an estimated model does not
repair an invalid observation assumption.

The financial interpretation follows the reward definition. Where costs are
already subtracted, subtracting contact cost again would double-count it.
Where a bank's reward law and initial distribution must also be estimated,
the supplied-law components of the generated comparison cannot be treated
as known. Outcome follow-up, action support and measured decision history
therefore determine whether its narrower uncertainty can be useful.

\chapter{A learned graph can leave the pricing decision unresolved}
\label{rewrite:discovery}

Suppose $X$ measures a proposed rate change in percentage points, $M$ an
intermediate market quantity and $Y$ gross funding contribution in thousands
of pounds. The units and coefficients in this section are authored for the
example. Consider the centered Gaussian observational covariance
\[
 \Sigma=\begin{pmatrix}
 1&0.6&0.48\\0.6&1.36&1.088\\0.48&1.088&1.8704
 \end{pmatrix}.
\]
The law satisfies $X\perp Y\mid M$. A chain $X\to M\to Y$, a
fork $X\leftarrow M\to Y$, and a reverse chain $X\leftarrow M\leftarrow Y$
can all reproduce it using independent Gaussian errors.

For the first chain, take $M=0.6X+\varepsilon_M$ and
$Y=0.8M+\varepsilon_Y$, with variances $\Var(X)=1$,
$\Var(\varepsilon_M)=1$ and $\Var(\varepsilon_Y)=1$. The intervention
slope of $Y$ on $X$ is $0.6(0.8)=0.48$. For the fork, take
$\Var(M)=1.36$, $X=(0.6/1.36)M+\eta_X$, and $Y=0.8M+\eta_Y$,
with residual variances $1-0.6^2/1.36$ and $1$. The resulting covariance
is the same, but intervening on $X$ has zero effect on $Y$. The reverse
chain uses the corresponding conditional regression coefficients and also
has zero effect. The numerical check reconstructs all three covariance
matrices and verifies positive residual variances.

\section{Keep the causal ambiguity in the financial report}
Within this specified three-model family, the slope set is $\{0,0.48\}$.
At a quarter-point intervention and a stipulated \pounds60 action cost,
net value is
\[
 1000(0.25)\{0,0.48\}-60=\{-60,60\}
 \quad\hbox{pounds}.
\]
There is no sign-robust recommendation across the retained explanations.
Returning the slope from the forward chain alone would hide the unresolved
orientation that determines the action. A probability for each graph would
require additional justified modelling; equal probabilities are not supplied
by observational equivalence. Nor should the discrete set be called a
confidence interval. It represents identification ambiguity even if the
covariance is known exactly.

The example's finite graph family and linear Gaussian law are restrictions.
They do not describe every model compatible with arbitrary banking records.
When the covariance and graph class are estimated, uncertainty must cover
both the observed law and the retained models before mapping them to effects.
An edge bootstrap conditional on one chosen graph does not supply that
combined guarantee. The foundations and detailed temporal-discovery
comparisons retain the arguments for unresolved orientations, latent causes
and class-to-effect propagation.

For a bank, useful additional information could be a defensible ordering
that excludes two orientations, a valid instrument, or an intervention on
$X$. The cost of obtaining that information should be compared with the
consequence of a mistaken action. More precise estimation of this same
covariance cannot distinguish these three laws. The proposed service should
therefore preserve a class or explicit refusal and ask what would narrow it,
instead of completing arrows solely to make a diagram look decisive.

\section{From a known covariance to estimated effect sets}
Chapter~\ref{exec13:temporal-region} supplies the next inferential step under
a restricted stationary Gaussian time-series model. It constructs a region
jointly over the observed covariance, a continuous unknown common lag and a
declared family of causal models, then maps that region to effects. The
coverage argument keeps the true model and observed-law uncertainty together;
it does not condition on an orientation chosen for convenience. The chapter
also shows what a direct intervention can add.

This is a developed restricted construction, not an executed general
bank-discovery pipeline. The maintained graph family, scalar lag and
innovation law must be defensible before applying it. An unrestricted latent
model may leave the whole line admissible. In the opening pricing example,
the narrower set already contains opposite net values, so a more elaborate
estimator is useful only if its information or restrictions can resolve a
consequential ambiguity. A labelled whole-line result can be the correct
answer to a well-defined question. A separate computational issue arises
when a search has not tested the conditioning sets needed to remove a link.
Section~\ref{repair29:discovery} compares a repair on exact-oracle cases and
shows why an unfinished search needs an explicit status. Search completeness,
statistical uncertainty and identification ambiguity are different limitations.

\chapter{Correcting records can change the action}
\label{rewrite:observation}

A treatment record is a measurement of an event. Let $A$ be true randomized
receipt, $W$ recorded receipt and $Y$ gross contribution in pounds at a
fixed horizon. In this separate finite example, $P(A=0)=P(A=1)=1/2$
and $\E[Y\mid A=0]=100$, $\E[Y\mid A=1]=140$. Receipt is
randomized by construction; correcting its record would not establish
randomization in an observational campaign.

Suppose misclassification is nondifferential with a known matrix
\[
 K=\begin{pmatrix}0.9&0.2\\0.1&0.8\end{pmatrix},
 \qquad K_{wa}=P(W=w\mid A=a,Y).
\]
Writing $p_a=P(A=a)$ and $q_a=\E[Y\mathbf1(A=a)]$, the recorded
shares and outcome moments obey
\[
 p_W=Kp_A=(0.55,0.45)' ,\qquad q_W=Kq_A=(59,61)' .
\]
Thus the recorded contrast is $61/0.45-59/0.55=28.2828\ldots$.
The matrix has determinant $0.7$, so its inverse recovers
$p_A=(0.5,0.5)'$ and $q_A=(50,70)'$, giving the true contrast
$70/0.5-50/0.5=40$. This is the moment version of the joint-law
argument in Chapter~\ref{r3:observation}. It is checked by direct
multiplication and solution, with no estimated audit law.

\section{The value of knowing the observation law}
At a stipulated action cost of \pounds35, the recorded comparison gives
net value about $-\pounds6.72$, whereas the corrected contrast gives
\pounds5. That reversal is conditional on the known error law. If an
identity matrix and the displayed matrix both remain admissible, the same
recorded shares and moments are compatible with different latent effects
and opposite recommendations. Reporting only the corrected number would
replace measurement uncertainty with a convenient assumption.

In a real validation sample, estimating $K$ would add uncertainty. The audit
must identify true receipt and represent the relevant customers and periods;
a classifier's own label is not an external truth measurement. Near-singular
matrices amplify error, and a singular matrix destroys recoverable information.
The retained chapter gives the symmetric-error example where receipt becomes
uninformative. More rows with the same deficient measurement law do not
restore its rank.

Entity linkage poses a related but distinct problem. A household may hold
several accounts, and a mistaken join can change both its treatment and its
outcome. Correctness counts alone may fail to identify the needed contrast;
the detailed linkage studies and their adverse results remain in this book.
The next useful observation is a correctly sampled, truthful linked history
when that is what the estimand requires, not merely a higher aggregate match
score. Likewise, event time, availability time and revision time must remain
separate when constructing a prospective decision view.

The four journeys now imply a common requirement for the proposed service.
It must keep the economic target and the information that identifies it
attached to the numerical result. Their statistical procedures remain
different: stratified assignment contrasts, sequential policy averaging,
model-set propagation and observation-law inversion. The next chapters
develop those differences before the software and investment synthesis chooses
which ones are ready for a bounded comparison.

\section{A fitted linkage analysis needs an external information budget}
Chapter~\ref{gap29:observation} studies a different observation problem:
matching randomized participants to outcomes when identifiers are corrupted.
It compares exact and fuzzy rules, keeps unmatched customers in the target,
and derives a contribution bound when at most a supplied number of candidate
lists omit the true match. Sampling uncertainty is then added under explicit
support and assignment assumptions. This is not an estimate of the receipt
matrix used above; the two studies concern different measurement mechanisms.

At 30\% outcome-dependent corruption, the robust population interval averages
119.1 outcome units in width, compared with about 4.1 for the fuzzy analysis.
The extra width prices justified ignorance. The corruption budget and outcome
support are external premises, not quantities learned from matching confidence.
A monetary decision first maps those outcome units and bounds into its own
expense equation. When the admissible set straddles that decision threshold,
an independently resolved audit sample can be more valuable than a forced
match. The detailed history and audit studies retain the further complications
of dependent records and uncertain audit truth.
